\documentclass[11pt]{scrartcl}

\usepackage[stixtwo]{fontsetup}
\usepackage{amsmath}
\usepackage{tikz}

\begin{document}
\begin{enumerate}
	\item \[
		      z
		      = \frac{(3 + 4i) - (2 - i)}{1 - i}
		      = \frac{1 + 5i}{1 - i}
		      = \frac{(1 + 5i)(1 + i)}{(1 - i)(1 + i)}
		      = \frac{-4 + 6i}{2}
		      = -2 + 3i,
	      \]
	      tehát $\operatorname{Re} z = -2$ és $\operatorname{Im} z = 3$.

	      Trigonometrikus alakban számolva
	      \(
	      \overline{2(\cos 15^\circ + i \sin 15^\circ)}
	      = 2(\cos(-15^\circ) + i \sin(-15^\circ))
	      \)
	      és $3 - 3i = 3\sqrt{2}(\cos(-45^\circ) + i \sin(-45^\circ))$, így
	      \[
		      w
		      = 6\sqrt{2}(\cos(-60^\circ) + i \sin(-60^\circ))
		      = 6\sqrt{2} \left( \frac{1}{2} - \frac{\sqrt{3}}{2} i \right)
		      = 3\sqrt{2} - 3\sqrt{6} i,
	      \]
	      tehát $\operatorname{Re} w = 3\sqrt{2}$ és
	      $\operatorname{Im} w = -3\sqrt{6}$. A szorzásnál a hosszak
	      összeszorzódnak, a szögek összeadódnak ($-15^\circ - 45^\circ = -60^\circ$):
	      \begin{center}
	      	\begin{tikzpicture}[scale=0.45, >=stealth]
	      		\draw[->] (-2.8, 0) -- (5.6, 0) node[right] {$\operatorname{Re}$};
	      		\draw[->] (0, -8) -- (0, 3.8) node[above] {$\operatorname{Im}$};
	      		\draw (0, 0) -- (3, -3);
	      		\draw (0, 0) -- (-15:2);
	      		\draw[thick] (0, 0) -- (-60:8.485);
	      		\fill (-2, 3) circle (2pt) node[above left] {$z$};
	      		\fill (3, -3) circle (2pt) node[right] {$3 - 3i$};
	      		\fill (-15:2) circle (2pt) node[below right] {$2 e^{-i 15^\circ}$};
	      		\fill (-60:8.485) circle (2pt) node[right] {$w$};
	      	\end{tikzpicture}
	      \end{center}

	\item $z^2 + 2z + 3i(z + 1) = z^2 + (2 + 3i) z + 3i$, így a gyökök
	      \[
		      z_{1, 2}
		      = \frac{-(2 + 3i) \pm \sqrt{(2 + 3i)^2 - 12i}}{2}
		      = \frac{-2 - 3i \pm \sqrt{-5}}{2}
		      = \frac{-2 - 3i \pm \sqrt{5} i}{2},
	      \]
	      azaz
	      \(
	      z_1 = -1 + \frac{\sqrt{5} - 3}{2} i
	      \)
	      és
	      \(
	      z_2 = -1 - \frac{\sqrt{5} + 3}{2} i
	      \).
	      A gyökök a komplex számsíkon:
	      \begin{center}
	      	\begin{tikzpicture}[scale=0.8, >=stealth]
	      		\draw[->] (-2.4, 0) -- (1.2, 0) node[right] {$\operatorname{Re}$};
	      		\draw[->] (0, -3.2) -- (0, 1.2) node[above] {$\operatorname{Im}$};
	      		\fill (-1, -0.382) circle (2pt) node[left] {$z_1$};
	      		\fill (-1, -2.618) circle (2pt) node[left] {$z_2$};
	      	\end{tikzpicture}
	      \end{center}
\end{enumerate}
\end{document}
